\documentclass[12pt, titlepage]{article}
\usepackage{amsmath}
\usepackage{amsfonts}
\usepackage{amssymb}
\usepackage{fancyhdr}
\usepackage{cancel}
\usepackage[shortlabels]{enumitem}
\usepackage[bidi=basic]{babel}
\babelprovide[main,import,Alph=letters]{hebrew}
\babelfont[hebrew]{rm}{Hadasim CLM}
\babelfont[hebrew]{sf}{Miriam Mono CLM}
\usepackage{titlesec}
\title{קומבינטוריקה 1040286 \\ גליון 2}
\author{יונתן אבידור - 214269565\\רותם עשת - 212674683}
\titleformat{\section}{\Large\sffamily\bfseries}{\thesection}{1em}{}
\begin{document}
\maketitle
\pagestyle{fancy}
\fancyhead{}
\fancyfoot[L]{\text{214269565}}
\fancyfoot[R]{\text{212674683}}
\part*{שאלה 1}
מצאו את מספר הפתרונות בשלמים אי שליליים למשוואות הבאות:
\section*{סעיף א'}
$2x_1+2x_2+x_3=11$ כאשר $x_2 \geq 1$
\section*{סעיף ב'}
$\left(x_1 + x_2 + x_3\right)\left(x_4+x_5\right)=18$
\part*{פתרון 1}
\section*{סעיף א'}
נשים לב כי $2x_1$ זוגי אי-שלילי ולכן קיים $a$ שלם אי-שלילי כך ש-$2x_1=2a$.\\
בנוסף לכך, $2x_2$ זוגי חיובי ולכן קיים $b$ שלם אי-שלילי כך ש-$2x_2 =2b+2$.\\
מכיוון ש-$2a+2b+2$ הוא מספר זוגי, ו-$11$ הוא מספר אי-זוגי, בהכרח $x_3$ הוא מספר אי-זוגי.\\
$x_3$ מספר אי-זוגי אי-שלילי ולכן קיים $c$ שלם אי-שלילי כך ש-$x_3 = 2c+1$\\
לכן
\begin{gather*}
    2x_1 + 2x_2 + x_3 = 11 \implies 2a + 2b + 2 + 2c  + 1 = 11 \\
    \implies 2\left(a+b+c\right) = 8 \implies a+b+c = 4
\end{gather*}
צמצנו את הבעיה למציאת מספר הפתרונות של משוואה בשלושה נעלמים שלמים אי-שליליים ללא תנאים נוספים ולכן כמות הפתרונות היא
\[
    \binom{4+3-1}{3-1} = \binom{6}{2} = \frac{6!}{2!4!} = \frac{720}{48} =\frac{60}{4} = \boxed{15}
\]
\section*{סעיף ב'}
נסמן
\begin{gather*}
    \mathbb{N} \ni  a := \left(x_1+x_2+x_3\right) \\
    \mathbb{N}  \ni b:= \left(x_4+x_5\right) \\
    a + b = 18 \implies b= \frac{18}{a}
\end{gather*}
נוכל לבחור את $a$ כמחלק כלשהו של $18$, ואז יש $\binom{a+2}{2}$ אפשרויות עבור $x_1,x_2,x_3$\\
לאחר מכן מכיוון ש-$b=\frac{18}{a}$, יהיו לנו $\binom{\frac{18}{a}+1}{1} =\frac{18}{a}+1$ אפשרויות עבור $x_4,x_5$ \\
על מנת למצוא את כמות האפשרויות עבור $x_{1-5}$, צריך לבחור בנפרד $x_{1-3}$ ואז עבור $x_{4-5}$.\\
לכן מכלל הכפל יש להכפיל את האפשרויות שלהם.\\ 
מכיוון ש-$a$ יכול להיות כל אחד מהמחלקים של $18$ בנפרד, מכלל החיבור נסכום את כמות האפשרויות עבור כל אחד מהם.\\
המחלקים של $18$ (האפשרויות של $a$) הם $1,2,3,6,9,18$ ולכן
\[
    \sum_{a\in \left\{1,2,3,6,9,18\right\}} \binom{a+2}{2} \cdot \left(\frac{18}{a} + 1\right)
\]
נחשב את הסכום ישירות
\begin{gather*}
    a=18,b=1 : \binom{20}{2} \cdot 2 = 190 \cdot 2 = 380 \\
    a=9,b=2 : \binom{11}{2} \cdot 3 = 55 \cdot 3 = 165 \\
    a=6,b=3 : \binom{8}{2} \cdot 4 = 28 \cdot 4 = 112 \\
    a=3,b=6 : \binom{5}{2} \cdot 7 = 10 \cdot 7 = 70 \\
    a=2,b=9 : \binom{4}{2} \cdot 10 = 6 \cdot 10  = 60 \\
    a=1,b=18: \binom{3}{2} \cdot 19 = 3 \cdot 19 = 57 \\
\boxed{380 + 165 + 112 + 70 + 60 + 57 = 844}
\end{gather*}
\newpage
\part*{שאלה 2}
כמה מספרים טבעיים בני $5$ ספרות יש שסכום ספרותיהם הוא 10?
\part*{פתרון 2}
נגדיר את ספרת האחדות $x_1$, העשרות $x_2$, המאות $x_3$, האלפים $x_4$\\
עבור ספרת העשרות-אלפים, מכיוון שהיא לא יכולה להיות $0$, נגדיר כ-$x_5+1$ \\
על מנת שסכום הספרות יהיה $10$, נדרוש שיתקיים:
\[
    x_1 + x_2 + x_3 + x_4 + x_5 + 1 = 10 \implies x_1 + x_2 + x_3 + x_4 + x_5 = 9
\]
קיבלנו משוואה של חמישה נעלמים שלמים אי-שליליים ולכן כמות האפשרויות היא
\[
    \binom{9+5-1}{5-1} = \binom{13}{4} = 715
\]
כל ספרה יכולה לתרום עד $9$ לסכום, אולם ספרת העשרות-האלפים היא $x_5+1$ \\
מכאן ש-$x_5$ עצמו יכול לתרום רק עד $8$ לסכום, לכן נרצה להוריד את כל האפשרויות בהן $x_5=9$\\
אך קיימת רק אפשרות אחת כזו:
\[
    x_1=0,x_2=0,x_3=0,x_4=0,x_5=9
\]
לכן כמות האפשרויות היא
\[
\boxed{715-1 = 714}
\]
\newpage
\part*{שאלה 3}
יהיו $n_1,n_2,n_3,k \in \mathbb{N}$. חשבו את הסכום:
\[
    \sum_{\begin{gathered}k_1+k_2+k_3=k\\k_i \geq 0 \end{gathered}} \binom{n_1}{k_1}\binom{n_2}{k_2}\binom{n_3}{k_3}
\]
\part*{פתרון 3}
נראה כי התשובה היא
\[
    \binom{n_1+n_2+n_3}{k}
\]
בכך שנראה כי שני הביטויים סופרים את אותו הדבר\\
תהא קבוצה בגודל $n_1+n_2+n_3$.\\
כמות האפשרויות לבחירה של $k$ איברים ממנה היא $\binom{n_1+n_2+n_3}{k}$\\
כעת נחלק את האיברים באופן שרירותי ל-$3$ קבוצות $n_1,n_2,n_3$ כך שגודלה של כל תת-קבוצה גדול או שווה ל-$k$ (זה אפשרי מכיוון שבניסוח המקורי של השאלה אנחנו עוברים על כל הצירופים של $k_1+k_2+k_3=k$, לכן זה כולל את הצירופים בהם שניים מה-$k_i$-ים הם $0$ והשלישי הוא $k$, אזי כל אחד מה$n_i$ גדול או שווה ל-$k$, אחרת לביטוי אין משמעות).\\
כאשר בוחרים $k$ איברים משלושת הקבוצות הללו יחדיו, קיימים $k_1,k_2,k_3$ כך שבוחרים $k_1$ איברים מתוך $n_1$, $k_2$ איברים מתוך $n_2$ ו-$k_3$ איברים מתוך $n_3$, וכן $k_1+k_2+k_3=k$.\\
עבור בחירה כלשהי של $k_1,k_2,k_3$, נכפיל את הבחירות של ה-$k_1$ מתוך ה-$n_1$ בבחירה של ה-$k_2$ מתוך ה-$n_2$ ושל ה-$k_3$ מתוך ה-$n_3$.\\
לאחר מכן, נסכום עבור האפשרויות לבחור את אותם $k_1,k_2,k_3$ ונקבל
\[
    \sum_{\begin{gathered}k_1+k_2+k_3=k\\k_i \geq 0 \end{gathered}} \binom{n_1}{k_1}\binom{n_2}{k_2}\binom{n_3}{k_3}
    \]
הראינו כי שני הדברים סופרים אותו דבר ולכן הביטויים שווים, ולכן
\[
    \boxed{\sum_{\begin{gathered}k_1+k_2+k_3=k\\k_i \geq 0 \end{gathered}} \binom{n_1}{k_1}\binom{n_2}{k_2}\binom{n_3}{k_3} = \binom{n_1+n_2+n_3}{k}}
\]
\newpage
\part*{שאלה 4}
חשבו את גדלי הקבוצות:
\section*{סעיף א'}
\[
    \left\{\left(A,B,C\right):A\subseteq B \subseteq C \subseteq \left[n\right] \right\}
\]
\section*{סעיף ב'}
\[
    \left\{\left(x_1,\ldots,x_n\right)\in\left\{0,1,2\right\}^{n}|\text{אפס מופיע מספר זוגי של פעמים }\right\}
\]
\part*{פתרון 4}
\section*{סעיף א'}
נגדיר פונקציה אשר מקבלת שלוש סדרות $A,B,C$ כנ"ל ומחזירה סדרה באורך $n$ מעל $\left\{1,2,3,4\right\}$ המוגדרת כך
\begin{gather*}
    \left(A,B,C\right) \mapsto \left(f_{A,B,C}\left(1\right),f_{A,B,C}\left(2\right),\ldots,f_{A,B,C}\left(n\right)\right)
\end{gather*}
נגדיר "עומק" של איבר $i\in\left[n\right]$ לפי מספר הקבוצות להן הוא שייך בשרשרת ההכלה של $A,B,C,\left[n\right]$\\
כלומר אם $i\notin C$ אזי העומק של $i$ הוא $1$ ואם $i\in A$ אזי העומק של $i$ הוא $4$\\
לפי זה, נגדיר את $f_{A,B,C}$
\begin{gather*}
    f_{A,B,C}:\left[n\right] \to \left\{1,2,3,4\right\} \\
    \forall i \in \left[n\right]: f_{A,B,C}\left(i\right) = \begin{cases}
        1 & i\in \left[n\right]\setminus C \\
        2 & i \in C \setminus B\\
        3 & i \in B \setminus A\\
        4 & i \in A\\
    \end{cases}
\end{gather*}
$f_{A,B,C}$ היא פונקציית העומק של הקבוצות $A,B,C$\\
נטען כי $F$ היא הפיכה, נעשה זאת בכך שנמצא את ההופכית שלה במפורש\\
נגדיר את הפונקציה הבאה
\begin{gather*}
    G: \left\{1,2,3,4\right\}^n \to \left(A',B',C'\right)
\end{gather*}
כאשר $A',B',C'$ קבוצות המקיימות $A' \subseteq B' \subseteq C' \subseteq \left[n\right]$\\
הפונקציה $G$ תפעל בדרך הבאה:\\
לכל $i\in \left[n\right]$, $G$ תכניס את $i$ לקבוצות המתאימות על פי $a_i$
\begin{gather*}
    a_i=1\implies i\in \left[n\right]\setminus C' \implies i\notin C' , i \notin B' , i \notin A'\\
a_i = 2 \implies i \in C' \setminus B' \implies i \in C' , i \notin B' , i \notin A' \\
a_i = 3 \implies i \in B' \setminus A' \implies i \in C', i \in B' , i \notin A'\\
a_i = 4 \implies i \in A' \implies  i \in B' , i \in C', i \in A'
\end{gather*}
לכן השלישייה הסדורה $\left(A',B',C'\right)$ אכן מקיימת את התנאי ש $A'\subseteq B'\subseteq C'\subseteq \left[n\right]$\\
נשים לב שעבור שלישייה סדורה המקיימת את התנאי $\left(\tilde{A},\tilde{B},\tilde{C}\right)$ מתקיים
\[
    G\left(F\left(\tilde{A},\tilde{B},\tilde{C}\right)\right)=\left(\tilde{A}, \tilde{B},\tilde{C}\right)
\]
שכן $F$ תיצור סדרה שבה כל $a_i$ שווה לעומק של $i$, ו-$G$ תהפוך את הסדרה הזו לשלוש קבוצות לפי העומק של כל איבר בסדרה\\
מאותה סיבה, עבור סדרה $a_n \in \left\{1,2,3,4\right\}^n$ מתקיים כי
\[
    F\left(G\left(a_n\right)\right) = a_n
\]
לכן $G=F^{-1}$, מכאן ש-$F$הפיכה, ולכן גודל התחום (כמות האפשרויות עבור $A,B,C$) שווה לגודל הטווח (כמות הסדרות באורך $n$ מעל $\{1,2,3,4\}$). לכן צריך עכשיו רק לחשב כמה סדרות שכאלה יש\\
יש לבחור $n$ איברים, כשהסדר משנה ומותר חזרות, מבין $4$ אפשרויות, לכן כמות האפשרויות היא
\[
    \boxed{4^n}
\]
\section*{סעיף ב'}
נקרא לסדרה בעלת כמות זוגית של אפסים "סדרה זוגית", ולסדרה בעלת כמות אי-זוגית של אפסים "סדרה אי-זוגית" \\
נראה באינדוקציה על $n$ כי כמות הסדרות הזוגיות באורך $n$ גדולה ב-$1$ מכמות הסדרות האי-זוגיות באורך $n$
\subsection*{בסיס האינדוקציה}
עבור $n=1$ מתקבלות הסדרות הבאות:\\
$\left(0\right)$ (סדרה אי-זוגית) \\
$\left(1\right)$ (סדרה זוגית) \\
$\left(2\right)$ (סדרה זוגית) \\
קיבלנו שיש $2$ סדרות זוגיות וסדרה זוגית אחת
\subsection*{הנחת האינדוקציה}
נניח שעבור סדרה באורך $n$ קיימות $k$ סדרות אי-זוגיות ו-$k+1$ סדרות זוגיות.
\subsection*{צעד האינדוקציה}
כל סדרה באורך $n+1$ מתקבלת על ידי הוספה של איבר כלשהו מבין $\left\{0,1,2\right\}$ לסוף של סדרה קיימת באורך $n$\\
עבור הוספה של האיבר $0$ לסוף סדרה באורך $n$, מתקבלות $k$ סדרות זוגיות ו-$k+1$ סדרות אי-זוגיות (מהנחת האינדוקציה והעובדה שהוספת $0$ משנה את הזוגיות של סדרה)\\
עבור הוספה של האיברים $1$ או $2$ לסוף סדרה באורך $n$, מתקבלות $k+1$ סדרות זוגיות ו-$k$ סדרות אי-זוגיות (מהנחת האינדוקציה והעובדה שלהוסיף $1$ או $2$ לא משנה את הזוגיות של הסדרה). אלו כל הדרכים לבנות סדרה באורך $n+1$ מעל $\left\{0,1,2\right\}$\\
נסכום את כמות הסדרות הזוגיות באורך $n+1$
\[
    \left(k\right) + \left(k+1\right) + \left(k+1\right) = 3k+2
\]
נסכום את כמות הסדרות האי-זוגיות באורך $n+1$
\[
    \left(k+1\right) + \left(k\right) + \left(k\right) = 3k+1
\]
ואכן קיבלנו כי כמות הסדרות הזוגיות גדולה ב-$1$ מכמות הסדרות האי-זוגיות כנדרש\\
כעת, מכיוון שכמות הסדרות באורך $n$ הוא $3^n$ (מספר אי זוגי). מתקבל כי כמות הסדרות הזוגיות באורך $n$ הוא
\[
    \boxed{  \lfloor \frac{3^n}{2}\rfloor+1 = \frac{3^n+1}{2}}
\]
\newpage
\part*{שאלה 5}
\section*{סעיף א'}
השתמשו בנוסחת הבינום כדי להוכיח:
\[
    \sum_{k=0}^{n} \binom{n}{k}^2 = \binom{2n}{n}
\]
\section*{סעיף ב'}
חשבו:
\[
    \sum_{k=0}^{n} \left(-1\right)^k \binom{n}{k}^2
\]
\part*{פתרון 5}
\section*{סעיף א'}
נתבונן בשוויון הבא המתקבל על פי חוקי חזקות
\[
    \left(x+1\right)^n \cdot \left(x+1\right)^n = \left(x+1\right)^{2n}
\]
נפתח את הביטויים בשני האגפים לפי נוסחת הבינום של ניוטון
\[
    \left(x+1\right)^n \cdot \left(x+1\right)^n = \sum_{k_1=0}^{n} \left(\binom{n}{k_1} \cdot x^{k_1}\right) \cdot \sum_{k_2=0}^{n}\left( \binom{n}{k_2} \cdot x^{k_2}\right)
\]
ומאגף ימין
\[
    \left(x+1\right)^{2n} = \sum_{k=0}^{2n} \left(\binom{2n}{k} \cdot x^{k}\right)
\]
מכיוון שהביטויים שווים, המקדמים של כל חזקה של $x$ שווים גם הם.\\
נשווה את המקדמים של $x^n$\\
מצד אחד, המקדם של $x^n$ הוא הסכום של כל הדרכים לבחור $n$ איברים משני הסכומים ביחד, כלומר סוכמים את מכפלת המקדמים לכל $k_1,k_2$ כך ש $k_1+k_2=n$, או אלגברית
\[
    \sum_{k=0}^{n} \left(\binom{n}{k}\cdot\binom{n}{n-k} \right) = \sum_{k=0}^{n}\left(\binom{n}{k}\cdot\binom{n}{k} \right) = \sum_{k=0}^{n} \binom{n}{k}^2\
\]
מהצד השני, המקדם של $x^n$ מתקבל רק כאשר $k=n$, ולכן הוא $\binom{2n}{n}$.
קיבלנו ש
\[
   \sum_{k=0}^{n} \binom{n}{k}^2 =   \binom{2n}{n}
\]
כנדרש\\
$\blacksquare$
\section*{סעיף ב'}
נתבונן בשוויון הבא המתקבל על פי חוקי חזקות
\[
    \left(1-x\right)^n \cdot \left(1+x\right)^n = \left(1-x^2\right)^n
\]
מצד שמאל נקבל
\[
    \sum_{k_1=0}^n \left(\binom{n}{k_1} \cdot \left(-1\right)^{k_1}\cdot x^{k_1}\right) \sum_{k_2=0}^n \left(\binom{n}{k_2}\cdot x^{n-k_2}\right)
\]
ומאגף ימין
\[
    \sum_{k=0}^{n} \left(\binom{n}{k} \cdot \left(-1\right)^{2k} \cdot x^{2k}\right)
\]
מכיוון שהביטויים שווים, המקדמים של כל חזקה של $x$ שווים גם הם.\\
נשווה את המקדמים של $x^n$\\
מצד אחד, המקדם של $x^n$ הוא הסכום של כל הדרכים לבחור $n$ איברים משני הסכומים ביחד, כלומר סוכמים את מכפלת המקדמים לכל $k_1,k_2$ כך ש
\[
    k_1+n-k_2=n\implies k_1=k_2
\]
או אלגברית
\[
    \sum_{k=0}^{n} \left(\left(-1\right)^{k} \binom{n}{k}^2 \right)
\]
מצד השני, $x^n$ מתקבל רק כאשר $k= \frac{n}{2}$, מה שקורה אם ורק אם $n$ זוגי, ואז הוא שווה ל$\binom{n}{\frac{n}{2}}\cdot\left(-1\right)^n$, ולכן
\[
\boxed{\sum_{k=0}^{n} \left(\left(-1\right)^{k} \binom{n}{k}^2 \right) = \begin{cases}
        \binom{n}{\frac{n}{2}} \cdot \left(-1\right)^n & n\text{ זוגי}\\
        0 & n \text{אי-זוגי }
\end{cases}}
\]
\newpage
\part*{בדיחת קרש}
למה שש פחד משבע\\\\\\
כי שבע שמונה תשע!
\end{document}
